\documentclass[12 pt, letterpaper]{hmcpset}
\usepackage{hw-preamble}

\name{Jayden Thadani}
\class{MATH 176 --- Algebraic Geometry}
\assignment{Homework assignment \#1}
\duedate{6th September 2024}

\begin{document}

\begin{problem}[1]
	If the sides of a triangle having a $60^{\circ}$ angle have sides of lengths $a$, $b$, and $c$, the law of cosines asserts that $a^{2} - ab + b^{2} = c^{2}$. We are interested in integer solutions $(a, b, c)$ to this equation. For any solution $(a, b, c)$, there are integers $m$ and $n$ such that $m/n = a/(c - b)$. Use this to show that
	\[
		(a : b : c) = (m^{2} + 2 m n : m^{2} - n^{2} : m^{2} + mn + n^{2}).
	\]
\end{problem}

\begin{solution}
	Write $t = m/n$. Then, substituting  $a = t(c - b)$. Then, since $(a, b, c)$ is a solution,
	\[
		t^{2}(c - b)^{2} - t(c - b)b + b^{2} = c^{2}.
	\]
	which we can rewrite as
	\[
		t^{2} (c - b)^{2} - t(c - b)b - (c - b)(c + b) = 0.
	\]
	If $c - b = 0$, the triangle must be an equilateral triangle, giving us $(a : b : c) = (1 : 1 : 1)$. In any other case, we must have
	\[
		(c - b) t^{2} - bt - (c + b) = 0.
	\]
	We can solve this for $b/c$.
	\[
		c t^{2} - c = b(t^{2} + t + 1)
	\]
	and thus,
	\[
		\begin{split}
			\frac{b}{c} & = \frac{t^{2} - 1}{t^{2} + t + 1} \\
				    & = \frac{m^{2} - n^{2}}{m^{2} + mn + n^{2}}.
		\end{split}
	\]
	Note then that we must have
	\[
		\begin{split}
			\frac{a}{b} & = \frac{t(c - b)}{b} \\
				    & = t \frac{c}{b} - t \\
				    & = \frac{t(t^{2} + t + 1) - t(t^{2} - 1)}{t^{2} - 1} \\
				    & = \frac{t^{2} + 2t}{t^{2} - 1} \\
				    & = \frac{m^{2} + 2mn}{m^{2} - n^{2}}.
		\end{split}
	\]
	With these two ratios, we have exactly the result we wanted ---
	\[
		(a : b : c) = (m^{2} + 2mn : m^{2} - n^{2} : m^{2} + mn + n^{2}).
	\]
\end{solution}

\pagebreak

\begin{problem}[2]
	Let $a$, $b$, $c$, and $d$ be rational numbers such that $ad$, $bd$, and $cd$ are not all zero (in particular $d$ is non-zero). Consider the intersection of the cone $z^{2} + y^{2} = z^{2}$ with a plane $ax + by + cz = d$.
	\begin{enumerate}
		\item Let $m$ and $n$ be integers. Show that
			\begin{align*}
				\frac{x}{z - y} & = \frac{m}{n} & & & x & = \frac{d (2 m n)}{a (2 m n) + b (m^{2} - n^{2}) + c(m^{2} + n^{2})} \\
				ax + by + cz & = d & \text{if and only if} & & y & = \frac{d (m^{2} - n^{2})}{a (2 m n) + b (m^{2} - n^{2}) + c(m^{2} + n^{2})} \\
				x^{2} + y^{2} & = z^{2} & & & z & = \frac{d (m^{2} + n^{2})}{a (2 m n) + b (m^{2} - n^{2}) + c(m^{2} + n^{2})}.
			\end{align*}
			Conclude that the intersection of the cone with the plane above has infinitely many rational points $(x, y, z)$.
		\item Show that this intersection is a conic section. That is, the curve
			\begin{align*}
				& & & & A & = c^{2} - a^{2} & D & = 2 a d \\
				A x^{2} + B x y + C y^{2} + D x + E y + F & = 0 & \text{in terms of} & & B & = -2 a b & E & = 2 b d \\
									  & & & & C & = c^{2} - b^{2} & F & = -d^{2}.
			\end{align*}
		\item Show that the equation $x^{2} + y^{2} + 1 = 0$ is not the intersection of a cone $x^{2} + y^{2} = z^{2}$ with a plane $ax + by + cz = d$ for any integers $a$, $b$, $c$, or $d$.
	\end{enumerate}
\end{problem}

\pagebreak

\begin{problem}[3]
	 Fix an integer $d$ which is not a square, and consider the equation $x^{2} - d y^{2} = 1$. Let $(z_{1}, y_{1})$ be an integral point on this curve, and define a sequence of integral points $\{ (x_{0}, y_{0}), (x_{1}, y_{1}), \dots, (x_{n}, y_{n}), \dots \}$ via the relation
	 \[
	 	x_{n} + y_{n} \sqrt{d} = (x_{1} + y_{1} \sqrt{d})^{n}
	 \]
	 \begin{enumerate}
		 \item By comparing the successive points $(x_{n}, y_{n})$ and $(x_{n - 1}, y_{n - 1})$, show that we have the matrix relation
			 \[
			 	\begin{bmatrix}
					x_{n} & d y_{n} \\
					y_{n} & x_{n}
			 	\end{bmatrix} = 
				\begin{bmatrix}
					x_{1} & d y_{1} \\
					y_{1} & x_{1} \\
				\end{bmatrix}
				\begin{bmatrix}
					x_{n - 1} & d y_{n - 1} \\
					y_{n - 1} & x_{n - 1}
				\end{bmatrix}.
			 \]
		\item Show that $x_{n}^{2} - d y_{n}^{2} = 1$.
	 \end{enumerate}
\end{problem}

\begin{solution}
	\begin{enumerate}
		\item We know that 
			\[
				\begin{split}
					x_{n} + y_{n} \sqrt{d} & =  (x_{1} + y_{1} \sqrt{d}) (x_{1} + y_{1} \sqrt{d})^{n - 1} \\
							       & = (x_{1} + y_{1} \sqrt{d})(x_{n - 1} + y_{n - 1} \sqrt{d}) \\
							       & = (x_{1} x_{n - 1} + d y_{1} y_{n - 1}) + (x_{1} y_{n - 1} + x_{n - 1} y_{1}) \sqrt{d}.
				\end{split}
			\]
			This gives us that $x_{n} = x_{1} x_{n - 1} + d y_{1} y_{n - 1}$ and $y_{n} = x_{1} y_{n - 1} + x_{n - 1} y_{1}$. Then
			\[
				\begin{split}
					\begin{bmatrix}
						x_{1} & d y_{1} \\
						y_{1} & x_{1}
					\end{bmatrix} 
					\begin{bmatrix}
					x_{n - 1} & d y_{n - 1} \\
					y_{n - 1} & x_{n - 1}
				\end{bmatrix} & = \begin{bmatrix}
				x_{1} x_{n - 1} + d y_{1} y_{n - 1} & d (x_{1} y_{n - 1} + x_{n - 1} y_{1}) \\
				x_{1} y_{n - 1} + x_{n - 1} y_{1} & x_{1} x_{n - 1} + d y_{1} y_{n - 1}
				\end{bmatrix} \\
				& = \begin{bmatrix}
					x_{n} & d y_{n} \\
					y_{n} & x_{n}
				\end{bmatrix}
				\end{split}
			\]
			as desired.
		\item
			Evidently, by induction
			\[
				\begin{bmatrix}
					x_{n} & d y_{n} \\
					y_{n} & x_{n}
				\end{bmatrix} = \begin{bmatrix}
					x_{1} & d y_{1} \\
					y_{1} & x_{1}
				\end{bmatrix}^{n}.
			\]
			Thus,
			\[
				\begin{split}
					x_{n}^{2} - d y^{2} & = \det \begin{bmatrix}
						x_{n} & d y_{n} \\
						y_{n} & x_{n}
					\end{bmatrix} \\
							    & = \left ( \det \begin{bmatrix}
						x_{1} & d y_{1} \\
						y_{1} & x_{1}
					\end{bmatrix} \right )^{n} \\
							    & = (x_{1}^{2} - d y_{1}^{2})^{n} \\
							    & = 1^{n} \\
							    & = 1.
				\end{split}
			\]
	\end{enumerate}
\end{solution}

\end{document}
